\title{FlashAttention-3: Fast and Accurate Attention with Asynchrony and Low-precision}

\begin{abstract}
Attention mechanisms form the foundational layer in the widely adopted Transformer model, acting as a computational bottleneck, particularly for large language models and scenarios demanding extended context. \textsc{FlashAttention} introduced a method to accelerate attention computations on GPUs by significantly reducing memory access operations. Nevertheless, prior solutions have not fully leveraged the latest enhancements available on modern hardware, as evidenced by \textsc{FlashAttention-2}, which utilizes only 35\% of the H100 GPU’s capacity. We introduce three key strategies to further accelerate attention on Hopper GPUs: utilizing the asynchronous operation of Tensor Cores and TMA to (1) concurrently execute data transfer and computation through warp specialization, (2) alternate between block-level matrix multiplications and softmax in a pipeline, and (3) utilize block quantization with non-coherent processing that makes use of hardware-based FP8 precision. Through these advances, our approach, \textsc{FlashAttention-3}, delivers a 1.5-2.0$\times$ performance improvement on H100 GPUs, reaching up to 740 TFLOPs/s in FP16 (75\% GPU utilization), and approaching 1.2 PFLOPs/s with FP8 precision. Our evaluation further shows that FP8 \textsc{FlashAttention-3} yields a 2.6$\times$ reduction in numerical error compared to a standard FP8 attention method.
\end{abstract}

\section{Introduction}
\label{sec:intro}

In the Transformer architecture~\citep{vaswani2017attention}, the principal computational constraint arises from the attention mechanism, as calculating self-attention scores between queries and keys results in a quadratic complexity with respect to the input sequence length. Enhancing the scalability of attention to encompass longer contexts is anticipated to facilitate advanced functionalities—such as handling and reasoning over numerous lengthy documents~\citep{guo2021longt5,shaham2022scrolls,peng2023yarn} and extensive source code repositories~\citep{roziere2023code, li2023starcoder}—as well as new data types, including detailed visual data~\citep{chen2022scaling}, audio signals~\citep{gulati2020conformer}, and video streams~\citep{ho2022video}. It also paves the way for innovative applications, including modeling user behavior across long temporal spans~\citep{sun2019bert4rec} and supporting complex agent workflows with prolonged planning horizons~\citep{yao2022react}. Consequently, there has been substantial research devoted to accelerating attention for long-context scenarios, such as through various approximation methods~\citep{katharopoulos2020transformers,choromanski2020rethinking,tay2020efficient}, advancements in software efficiency~\citep{rabe2021self,dao2022flashattention,kwon2023efficient}, and even the exploration of novel model architectures~\citep{peng2023rwkv,sun2023retentive,gu2023mamba}.

Our study extends the contributions of \citet{dao2022flashattention}, who pioneered exact-attention algorithms by embedding detailed knowledge of the GPU execution pipeline and underlying hardware features into algorithmic design. In their foundational work, Dao et al. proposed \textsc{FlashAttention}, introducing an innovative tiling approach to parallelize attention computations; by consolidating all attention operations within a single GPU kernel, their method avoided the costly overhead of intermediate accesses to slower global memory.

Later, \citet{dao2023flashattention2} revised this methodology with \textsc{FlashAttention-2}, restructuring the computation to allow parallelization over the sequence length as well. This refinement involved processing segments of the key and value matrices within the forward pass’s inner loop, thereby enhancing both work allocation and resource usage on the GPU. Despite these advancements, our analysis reveals that \textsc{FlashAttention-2} still underperforms, exhibiting substantially lower utilization rates on contemporary GPUs compared to highly-optimized matrix-multiplication (GEMM) routines—achieving only around 35\% utilization relative to 80-90\% for GEMMs on Hopper H100 GPUs. Such inefficiencies may, in part, result from a lack of tailoring the implementation to Hopper-specific architectural features, for example by persisting in using instruction sets optimized for Ampere rather than Hopper when deploying computations on Tensor Cores. Recent works—including ThunkerKitten~\citep{spector2024thunder} and cuDNN 9~\citep{cudnn9}—demonstrate that leveraging Hopper-native instructions and adopting tile-oriented computational strategies both accelerates attention mechanisms and streamlines their programming.

At a deeper level, the algorithm underlying \textsc{FlashAttention-2} is based on a straightforward, synchronous execution paradigm and does not incorporate asynchrony or reduced-precision computation within its core methodology. The emergence of asynchrony is primarily a consequence of recent advances in hardware specialization tailored to accelerate core machine learning operations: for instance, distinct components such as Tensor Cores are engineered for matrix multiplication while dedicated modules like the Tensor Memory Accelerator (TMA) handle memory operations. These units operate independently from the remaining CUDA cores, which are responsible for logical, integer, and floating-point calculations. Employing lower-precision formats—for example, FP8 in Hopper and FP4 in Blackwell, continuing the shift initiated with FP16 (introduced in Pascal, 2017) and BF16 (featured in Ampere, 2020)—enables significantly greater throughput without additional costs in power consumption or silicon real estate. We further explore the enhancements made available by Hopper in these respects in \cref{subsec:hardware}. The main technical obstacle involves adapting the \textsc{FlashAttention-2} algorithm to exploit such architectural capabilities: leveraging asynchrony necessitates finding ways to interleave the execution of matmul and softmax, despite inherent data dependencies between them; meanwhile, effective use of low-precision computation calls for strategies to limit quantization errors, a challenge heightened by the presence of extreme-valued features in LLMs~\citep{dettmers2208llm, sun2024massive}.

With this goal in mind, we introduce \textsc{FlashAttention-3}, which integrates and advances three novel concepts aimed at enhancing performance on cutting-edge GPU architectures:\footnote{Our results are presented with reference to NVIDIA's Hopper architecture, though our approach is applicable to any GPU architecture equipped with strong asynchronous execution and support for low-precision operations.}

\begin{enumerate}

\item \textbf{Producer-Consumer asynchrony:} We introduce a warp-specialized software pipelining approach that leverages asynchronous processing between data transfer and Tensor Core computations. By allocating producer and consumer tasks to different warps, the scheme expands the algorithm's potential to mask both memory and instruction issue latency.
\item \textbf{Hiding softmax under asynchronous block-wise GEMMs:} The algorithm overlaps the execution of the relatively slower non-GEMM softmax-related operations—such as floating point multiply-add and exponentiation—with the asynchronous WGMMA-based GEMM computations. In doing so, the \textsc{FlashAttention-2} algorithm is refactored to remove certain serialized dependencies between the softmax stage and GEMM operations. For instance, in our 2-stage approach, while one block of the score matrix undergoes softmax, the asynchronous WGMMA proxy is used to process the next matrix block in parallel.
\item \textbf{Hardware-accelerated low-precision GEMM:} To achieve nearly a two-fold increase in observed TFLOPs/s, we modify the forward pass algorithm to utilize FP8 Tensor Cores for GEMM. This adaptation involves reconciling the memory layout constraints imposed by WGMMA when handling FP32 accumulator and FP8 operand matrices. Block quantization and incoherent processing strategies are employed to alleviate the loss of numerical fidelity caused by the transition to FP8 representation.
\end{enumerate}

In order to empirically substantiate our approach, we conduct benchmarking of \textsc{FlashAttention-3} on the H100 SXM5 GPU, examining its performance across diverse parameter settings. The results demonstrate that (1) in the forward pass, the FP16 variant provides a 1.5-2.0$\times$ increase in speed relative to \textsc{FlashAttention-2}, achieving peak throughput of 740 TFLOPs/s, and offers a 1.5-1.75$\times$ improvement during the backward pass; (2) the FP8 mode attains nearly 1.2 PFLOPs/s; and (3) with extended sequence lengths, FP16 delivers superior performance and FP8 remains competitive\footnote{To specify, for head dimension 64, \textsc{FlashAttention-3} FP8 is ahead, whereas for head dimensions of 128 and 256, it matches performance in the absence of causal masking but falls behind when causal masking is present.} compared with the leading attention implementation available in NVIDIA’s cuDNN library. Our experiments further confirm that the FP16 configuration of \textsc{FlashAttention-3} achieves numerical errors indistinguishable from those of \textsc{FlashAttention-2}, outperforming conventional attention implementations, as intermediate computations—such as softmax rescaling—are retained in FP32. Additionally, FP8 \textsc{FlashAttention-3}, utilizing block quantization combined with incoherent processing, attains a 2.6$\times$ enhancement in accuracy over standard attention approaches that rely on per-tensor quantization, particularly in scenarios involving outlier features.

\textsc{FlashAttention-3} is released as open-source software under a permissive license\footnote{\textsc{FlashAttention-3}
  is available at \texttt{https://github.com/Dao-AILab/flash-attention}}. We also intend to incorporate support for it in both PyTorch and Hugging Face libraries, aiming to maximize its accessibility and utility for the research and developer community.

\section{Background: Multi-Head Attention and GPU Characteristics}
\label{sec:background}

\subsection{Multi-Head Attention}
\label{subsec:multi_head_attn}

Let $\mathbf{Q}, \mathbf{K}, \mathbf{V} \in \mathbb{R}^{N \times d}$ be the query, key and value input sequences associated to a single head, where $N$ is the sequence length and $d$ is the head dimension. Then the attention output $\mathbf{O}$ is computed as:

\begin{equation*}
  \mathbf{S} = \alpha \mathbf{Q} \mathbf{K}^\top \in \mathbb{R}^{N \times N}, \quad \mathbf{P} = \mathrm{softmax}(\mathbf{S}) \in \mathbb{R}^{N \times N}, \quad \mathbf{O} = \mathbf{P}\mathbf{V} \in \mathbb{R}^{N \times d},
\end{equation*}

Here, $\mathrm{softmax}$ operates on each row, and the scaling parameter $\alpha$ is conventionally chosen as $1/\sqrt{d}$. To mitigate numerical instability arising from exponentiation, it is common to subtract $\mathrm{rowmax}(\mathbf{S})$ from $\mathbf{S}$ before applying the softmax. In the case of multi-head attention (MHA), distinct sets of query, key, and value projections are used for each head. This allows the computation to be executed in parallel across several heads and batches, resulting in the complete output tensor.

Now let $\phi$ be a scalar loss function and let $\mathbf{d}(-) = \partial \phi / \partial (-)$ be notation for the gradient. Given the output gradient $\mathbf{dO} \in \mathbb{R}^{N \times d}$, we compute $\mathbf{dQ}$, $\mathbf{dK}$, and $\mathbf{dV}$ according to the chain rule as follows:

\begin{align*}
  \mathbf{dV} &= \mathbf{P}^\top \mathbf{dO} \in \mathbb{R}^{N \times d} \\
  \mathbf{dP} &= \mathbf{dO} \mathbf{V}^\top \in \mathbb{R}^{N \times N} \\
  \mathbf{dS} &= \mathrm{dsoftmax} (\mathbf{dP}) \in \mathbb{R}^{N \times N} \\
  \mathbf{dQ} &= \alpha \mathbf{dS} \mathbf{K} \in \mathbb{R}^{N \times d} \\
  \mathbf{dK} &= \alpha \mathbf{dS}^\top \mathbf{Q} \in \mathbb{R}^{N \times d},
\end{align*}

In this context, $\mathbf{d}s = (\mathrm{diag}(p) - p p^\top)\mathbf{d}p$, where $p = \mathrm{softmax}(s)$, expresses the relationship as a function of the vector $s$. We denote applying this formula to each row as $\mathrm{dsoftmax}(\mathbf{dP})$. As with previous operations, this procedure is efficiently parallelized over heads and batch elements during the MHA backward pass.

\subsection{GPU hardware characteristics and execution model}
\label{subsec:hardware}

We outline elements of the GPU execution model pertinent to \textsc{FlashAttention-3}, specifically emphasizing the NVIDIA Hopper architecture as a representative example of this framework.

\paragraph{Memory hierarchy:} The memory architecture of a GPU is structured as a hierarchy of data storage regions, with each successive level exhibiting higher bandwidth but lower capacity (\cref{tab:gpu-hierarchy})\footnote{\citet{luo2024benchmarking} measures shared memory bandwidth at 128 bytes per clock cycle per SM; this value is then scaled by 132 SMs and a boost clock frequency of 1830 MHz.}. At the top of this hierarchy is global memory (GMEM), or HBM, which is the off-chip DRAM available to all streaming multiprocessors (SMs). When data is accessed from GMEM, it is automatically cached in the on-chip L2 cache. Further down, every SM is equipped with its own fast, highly banked on-chip cache called shared memory (SMEM), which must be explicitly managed by programmers. The lowest level consists of the register file housed within each SM.

\paragraph{Thread hierarchy:} Within the GPU programming paradigm, execution units are structured into distinct logical groups referred to as threads. Arranged from the most granular to the broadest scale, the thread hierarchy consists of individual threads, warps (each comprising 32 threads), warpgroups (made up of 4 sequential warps), threadblocks (also known as cooperative thread arrays or CTAs), threadblock clusters (as introduced in Hopper), and grids.

The two hierarchies are tightly connected. All threads within a single CTA are simultaneously scheduled onto the same SM, while CTAs grouped in a cluster are jointly scheduled onto a single GPC. Every thread in a CTA can directly access SMEM, but each individual thread is limited to a private allocation of up to 256 registers (RMEM).

\paragraph{Asynchrony and warp-specialization:}

GPUs function as throughput-oriented processors, leveraging parallelism and asynchronous operations to mask both memory and execution delays. In the Hopper architecture, the Tensor Memory Accelerator (TMA) is introduced as specialized hardware to handle asynchronous memory transfers between GMEM and SMEM \cite[\S7.29]{cuda}. In addition, diverging from earlier architectures like Ampere, Hopper's Tensor Core—utilized through the warpgroup-level WGMMA instruction \cite[\S9.7.14]{ptx}—operates asynchronously and is capable of directly accessing its input data from shared memory.

By facilitating asynchrony at the hardware level, support is provided for warp-specialized kernels. In these kernels, a CTA’s warps are partitioned into producer and consumer groups—producers handle only data movement, while consumers focus solely on computation. This division generally enhances the compiler's capacity to emit more efficient instruction schedules \citep{warp-specialization-2011}. Furthermore, Hopper introduces the ability to dynamically redistribute registers across warpgroups through \verb|setmaxnreg| \cite[\S9.7.17.1]{ptx}. Consequently, warps executing MMAs can be allocated a greater portion of RMEM resources compared to those solely engaged in TMA operations, which require just a single participating thread.

\paragraph{Low-precision number formats:}
\label{sec:low-precision-gpu}
Recent generations of GPUs are equipped with dedicated components that enhance the processing speed of low-precision arithmetic. As an illustration, the WGMMA operation is capable of utilizing the FP8 Tensor Cores found in Hopper architectures, achieving double the throughput per SM relative to FP16 or BF16 formats.

Nevertheless, to utilize FP8 WGMMA correctly, one must account for specific operand layout requirements. Consider a GEMM operation computing $A \times B^{\top}$, where $A$ is an $M\times K$ matrix and $B$ is an $N\times K$ matrix. An operand is described as \emph{mn-major} when the data are stored contiguously along the outer $M$ (for $A$) or $N$ (for $B$) dimension; it is labeled \emph{k-major} if contiguousness occurs along the $K$ dimension. For FP16 WGMMA, shared memory (SMEM) operands may be formatted in either mn-major or k-major order. By contrast, FP8 WGMMA imposes stricter requirements, permitting only the k-major layout for input operands. Additionally, in scenarios like attention where chaining multiple GEMMs within one kernel is desirable, conflicting memory layouts between FP32 accumulators and FP8 operands can impede the invocation of successive FP8 WGMMA operations.

With regard to attention mechanisms, these layout constraints necessitate specific alterations to how an FP8 algorithm is constructed, as outlined in \cref{sec:algofp8}.

\subsection{Standard Attention and Flash Attention} Building on the definitions from \citet{dao2022flashattention}, we use the term \textbf{standard attention} to refer to an attention computation on GPU where the intermediate matrices $\mathbf{S}$ and $\mathbf{P}$ are explicitly stored in HBM. In contrast, the central innovation behind \textsc{FlashAttention} involves the application of a local softmax reduction, which significantly reduces the need for costly intermediate memory transactions by consolidating the attention operation into a single GPU kernel. This local softmax process maps to lines \ref{code-ws:softmax_start}-\ref{code-ws:softmax_end} within the consumer mainloop of \cref{alg:flash3_wgmma_ws_only}, coupled with the necessary block-wise scaling of $\mathbf{O}$. For a straightforward derivation demonstrating that this approach correctly computes $\mathbf{O}$, consult \cite[\S 2.3.1]{dao2023flashattention2}.

\section{FlashAttention-3: Algorithm}
\label{sec:algo}

In this section, we present an overview of the \textsc{FlashAttention-3} algorithm. Our primary emphasis is on the forward pass, while the procedure for the backward pass is provided in~\cref{sec:algo_ws_bwd}. Initially, we outline the incorporation of warp-specialization alongside a circular SMEM buffer into the original framework of \textsc{FlashAttention-2}. Subsequently, we discuss how leveraging the asynchronicity of WGMMA enables the construction of a two-stage GEMM-softmax pipeline with overlapping execution. Lastly, we detail the changes required to support FP8, addressing both memory layout alignment and the preservation of accuracy through the application of block quantization and the management of incoherent computations.

\subsection{Producer-Consumer asynchrony through warp-specialization and pingpong scheduling}
\label{sec:algo_ws}

\paragraph{Warp-specialization}
Similar to \textsc{FlashAttention-2}, the forward computation in \textsc{FlashAttention-3} is highly parallelizable across batch dimension, head count, and length of the query sequence. Therefore, a CTA-level abstraction is sufficient to capture the algorithm’s operation: it processes a query tile $\mathbf{Q}_i$ from the query matrix and generates the associated output tile $\mathbf{O}_i$. For clarity, we first describe the warp-specialization approach assuming a circular SMEM buffer is used \emph{without} introducing additional GEMM-softmax overlap. If we denote the head size by $d$ and the sequence length by $N$, and set the query block size as $B_r$, then $\mathbf{Q}$ is partitioned into $T_r = \lceil \frac{N}{B_r} \rceil$ blocks: $\mathbf{Q}_1, .., \mathbf{Q}_{T_r}$.

\begin{algorithm}[H]
    \caption{\small\label{alg:flash3_wgmma_ws_only}\textsc{FlashAttention-3} forward pass \textbf{excluding} intra-consumer overlapping -- CTA perspective}
    \begin{algorithmic}[1]
\REQUIRE The matrices $\mathbf{Q}_i \in \mathbb{R}^{B_r \times d}$ along with $\mathbf{K}, \mathbf{V} \in \mathbb{R}^{N \times d}$ are present in HBM; set the key block size $B_c$, calculate $T_c = \lceil \frac{N}{B_c} \rceil$.
\STATE Set up a pipeline object responsible for synchronization across barriers, utilizing an $s$-stage circular SMEM buffer.
\IF {inside the producer warpgroup}
\STATE Free up a preassigned number of registers.
\STATE Trigger the transfer of $\mathbf{Q}_i$ from HBM into shared memory.
\STATE Once the transfer is complete, notify the consumer that $\mathbf{Q}_i$ is available.
\FOR{$0 \le j < T_c$}
    \STATE Wait until the $(j\,\%\,s)$ stage of the buffer has been processed by the consumer.
    \STATE Start loading $\mathbf{K}_j, \mathbf{V}_j$ from HBM to shared memory, targeting the $(j\,\%\,s)$ slot.
    \STATE Upon finishing the memory copy, signal the consumers that $\mathbf{K}_j, \mathbf{V}_j$ are now ready.
\ENDFOR
\ELSE \STATE Restore the predetermined number of registers, as determined by the consumer warp count.
\STATE Locally initialize $\mathbf{O}_i = (0) \in \mathbb{R}^{B_r \times d}$, together with $\ell_i, m_i = (0), (-\infty) \in \mathbb{R}^{B_r}$.
\STATE Await the completion of loading $\mathbf{Q}_i$ into shared memory.
\FOR{$0 \le j < T_c$}
\STATE Wait for $\mathbf{K}_j$ to be present in shared memory.
\STATE Perform the operation $\mathbf{S}_i^{(j)} = \mathbf{Q}_i \mathbf{K}_j^T$ via SS-GEMM. Commit and block until completion. \label{alg:ws_only_gemm1}
\STATE Save $m_i^{\mathrm{old}} = m_i$, then update $m_i = \mathrm{max}(m_i^{\mathrm{old}}, \mathrm{rowmax}(\mathbf{S}_i^{(j)}))$. \label{code-ws:softmax_start}
\STATE Compute $\widetilde{\mathbf{P}}_i^{(j)} = \mathrm{exp}(\mathbf{S}_i^{(j)} - m_i)$, and update $\ell_i = \mathrm{exp}(m_i^{\mathrm{old}} - m_i) \ell_i + \mathrm{rowsum}(\widetilde{\mathbf{P}}_i^{(j)})$. \label{code-ws:softmax_end}
\STATE Wait until $\mathbf{V}_j$ is available in shared memory.
\STATE Evaluate $\mathbf{O}_i = \mathrm{diag}(\mathrm{exp}(m_i^{\mathrm{old}} - m_i))^{-1} \mathbf{O}_i + \widetilde{\mathbf{P}}_i^{(j)} \mathbf{V}_j$ through RS-GEMM. Commit and block as needed. \label{alg:ws_only_gemm2}
\STATE Release the $(j\,\%\,s)$ buffer stage to signal the producer it is free.
\ENDFOR
\STATE Finalize by computing $\mathbf{O}_i = \mathrm{diag}(\ell_i)^{-1} \mathbf{O}_i$ and $L_i = m_i + \log(\ell_i)$.
\STATE Commit $\mathbf{O}_i$ and $L_i$ to HBM, populating the $i$th blocks of $\mathbf{O}$ and $L$, respectively.
\ENDIF
\end{algorithmic}
\end{algorithm}

In our implementation of \cref{alg:flash3_wgmma_ws_only} on Hopper, we utilize \verb|setmaxnreg| to handle register (de)allocation, employ TMA to load $\mathbf{Q}_i$ as well as $\{\mathbf{K}_j, \mathbf{V}_j\}_{0 \leq j < T_c}$, and apply WGMMA for executing GEMMs within the consumer mainloop; the SS and RS prefixes denote whether the initial operand comes from shared memory or the register file, respectively. To properly interpret the execution logic of \cref{alg:flash3_wgmma_ws_only}, it is important to recognize that TMA loads are asynchronous, so new loads can be issued without waiting for prior ones to finish. Additionally, during the initial $s$ iterations of the producer mainloop, waits are not triggered, as the buffer is being initially populated.

\paragraph{Pingpong scheduling} The asynchronous operations of WGMMA and TMA, in combination with warp-specialization, make it possible to concurrently execute the softmax calculations for one warpgroup while another warpgroup processes GEMM. This overlapping is advantageous given that, on current hardware accelerators, the throughput of non-matmul tasks is significantly lower than that of matmul. For instance, on the H100 SXM5 GPU, the FP16 matmul achieves a throughput of 989 TFLOPS, while special functions such as the exponential operation\footnote{According to the CUDA programming guide, each streaming multiprocessor (SM) can execute 16 special function operations per clock cycle. By multiplying 16 by 132 SMs and the clock frequency of 1830 MHz, we obtain 3.9 TFLOPS for special functions.} (required for softmax) reach only 3.9 TFLOPS. In the FP16 attention forward pass with a head dimension of 128, the required number of matmul FLOPS is 512 times that of exponentials, but since the latter runs at 256 times less throughput, exponentials can occupy up to 50\% of the compute cycles relative to matmul. This inefficiency is exacerbated in FP8, where the matmul speed doubles, yet exponential throughput does not increase.

The exponential operation is handled by a dedicated hardware component known as the multi-function unit. Optimally, it is beneficial to overlap the computation of the exponential function with the matrix multiplication tasks being executed on the Tensor Cores. To achieve this overlap, we employ synchronization barriers (\texttt{bar.sync} instructions), ensuring that the GEMMs (specifically, GEMM1—$\mathbf{P} \mathbf{V}$ from the current iteration—and GEMM0—$\mathbf{Q} \mathbf{K}^\top$ from the subsequent iteration) for warpgroup 1 are dispatched before those for warpgroup 2. Consequently, while warpgroup 1 is computing the softmax, warpgroup 2 undertakes its GEMM operations. Following this, the responsibilities are reversed: warpgroup 2 computes softmax as warpgroup 1 executes the next set of GEMMs, a scheme referred to as "pingpong" scheduling, which is depicted in~\cref{fig:pingpong_scheduling}. Although this pingpong execution pattern is typically not as well-defined in actual practice as shown in the diagram, this technique generally yields performance benefits (for example, increasing throughput from 570 TFLOPS to between 620 and 640 TFLOPS for FP16 forward passes with a head dimension of 128 and a sequence length of 8192).

\paragraph{Attention variants} In the case of multi-query attention~\citep{shazeer2019fast} and grouped query attention~\citep{ainslie2023gqa}, we adopt the methodology from \textsc{FlashAttention-2}, modifying the tensor indexing so that duplication of $\mathbf{K}$ and $\mathbf{V}$ within HBM is avoided.

\subsection{Intra-warpgroup overlapping GEMMs and softmax}

It is possible to interleave certain instructions from the softmax computation with those from the GEMMs within a single warpgroup. Here, we outline a method for achieving this overlap.

In the attention algorithm, operations within the inner loop (main loop) have sequential dependencies that impede parallelization within a single iteration. For example, (local) softmax (lines \ref{code-ws:softmax_start} to \ref{code-ws:softmax_end}) relies on the output $\mathbf{S}_i^{(j)}$ of the first GEMM, while the second GEMM takes its result $\widetilde{\mathbf{P}}_i^{(j)}$ as an operand. Indeed, the wait statements in lines \ref{alg:ws_only_gemm1} and \ref{alg:ws_only_gemm2} of \cref{alg:flash3_wgmma_ws_only} serialize the execution of softmax and GEMMs. However, we can break these dependencies by pipelining across iterations through additional buffers in registers. Pursuing this idea, we propose the following two-stage\footnote{Note that the number of stages of the overlapping scheme is bounded by, but need not equal, the number $s$ of stages in the circular SMEM buffer.} GEMM-softmax pipelining algorithm:

\begin{algorithm}[H]
    \caption{\small\label{alg:flash3_wgmma}\textsc{FlashAttention-3} consumer warpgroup forward pass}
    \begin{algorithmic}[1]
      \REQUIRE Matrices $\mathbf{Q}_i \in \mathbb{R}^{B_r \times d}$ and $\mathbf{K}, \mathbf{V} \in \mathbb{R}^{N \times d}$ reside in HBM; the key block size $B_c$ is given, and $T_c = \lceil \frac{N}{B_c} \rceil$.
      \STATE Adjust the allocation of registers according to the total number of consumer warps.
      \STATE Allocate on-chip storage for $\mathbf{O}_i = (0) \in \mathbb{R}^{B_r \times d}$ along with $\ell_i, m_i = (0), (-\infty) \in \mathbb{R}^{B_r}$.
      \STATE Synchronize until both $\mathbf{Q}_i$ and $\mathbf{K}_0$ are loaded into shared memory.
      \STATE Perform the computation $\mathbf{S}_{\mathrm{cur}} = \mathbf{Q}_i \mathbf{K}_0^T$ using WGMMA; ensure the result is committed, then wait for its completion.
      \STATE Free the first stage ($0$) of the buffer dedicated to $\mathbf{K}$.
      \STATE Derive $m_{i}$, $\tilde{\mathbf{P}}_{\mathrm{cur}}$, and $\ell_{i}$ from $\mathbf{S}_{\mathrm{cur}}$, and rescale $\mathbf{O}_i$ accordingly.
      \FOR{$1 \le j < T_c - 1$}
        \STATE Wait until $\mathbf{K}_j$ becomes available in shared memory.
        \label{code:mainloop_start}
        \STATE Execute $\mathbf{S}_{\mathrm{next}} = \mathbf{Q}_i \mathbf{K}_{j}^T$ using WGMMA; commit but postpone waiting for completion. \label{code:first_wgmma}
        \STATE Block until $\mathbf{V}_{j-1}$ has been loaded into shared memory.
        \STATE Perform the update $\mathbf{O}_{i} = \mathbf{O}_{i} + \tilde{\mathbf{P}}_{\mathrm{cur}} \mathbf{V}_{j-1}$ with WGMMA, commit without immediate synchronization. \label{code:second_wgmma}
        \STATE Await the result of $\mathbf{Q}_i \mathbf{K}_{j}^T$ from WGMMA.
        \STATE Based on $\mathbf{S}_{\mathrm{next}}$, calculate $m_{i}$, $\tilde{\mathbf{P}}_{\mathrm{next}}$, and $\ell_{i}$. \label{code:softmax}
        \STATE Ensure $\tilde{\mathbf{P}}_{\mathrm{cur}} \mathbf{V}_{j-1}$ via WGMMA has completed before applying any further rescaling to $\mathbf{O}_i$.
        \STATE Release the $(j\,\%\,s)$th stage of $\mathbf{K}$ and $(j-1\,\%\,s)$th stage of $\mathbf{V}$ buffer.
        \STATE Move contents of $\mathbf{S}_{\mathrm{next}}$ into $\mathbf{S}_{\mathrm{cur}}$.
        \label{code:mainloop_end}
      \ENDFOR \STATE Block until $\mathbf{V}_{T_c - 1}$ has finished loading into shared memory.
      \STATE Carry out
      $\mathbf{O}_{i} = \mathbf{O}_{i} + \tilde{\mathbf{P}}_{\mathrm{last}} \mathbf{V}_{T_c - 1}$ by means of WGMMA; commit and then synchronize completion.

\STATE Epilogue: Adjust $\mathbf{O}_{i}$ by scaling it according to $m_{i}$. Determine $L_{i}$ utilizing both $m_{i}$ and $\ell_{i}$.
      Store the updated $\mathbf{O}_{i}$ and $L_{i}$ in HBM as the $i$-th segment of $\mathbf{O}$ and $L$.
    \end{algorithmic}
  \end{algorithm}

\cref{alg:flash3_wgmma} substitutes for the consumer pathway outlined in \cref{alg:flash3_wgmma_ws_only}, thereby presenting the full \textsc{FlashAttention-3} algorithm targeting FP16 computations. Broadly speaking, WGMMA is utilized here as a stand-in for asynchronous GEMM. During execution of the mainloop (spanning lines \ref{code:mainloop_start} through \ref{code:mainloop_end}), the second invocation of WGMMA in the $j$th iteration (at line \ref{code:second_wgmma}) is scheduled to execute concurrently with the softmax computations that correspond to iteration $j+1$ (found at line \ref{code:softmax}).

Although the pipelined design discussed above promises theoretical improvements in performance, several real-world factors must be taken into account:
\paragraph{Compiler reordering} The pseudocode assumes an ideal, sequential instruction execution, but in practice, the NVCC compiler frequently reorders operations to enhance efficiency. This optimization process can interfere with the deliberate interleaving of WGMMA and non-WGMMA instructions, possibly resulting in unpredictable behavior or a reduction in anticipated speedup. Nevertheless, inspection of the generated SASS confirms that, in practice, code overlapping as described is indeed produced (see Section~\ref{sec:2-stage-sass}).
\paragraph{Register pressure} Achieving peak throughput necessitates avoiding register spilling. However, implementing a 2-stage pipeline introduces a requirement for supplementary registers to hold temporary data and sustain stage-to-stage information. In particular, the algorithm needs to keep an additional $\mathbf{S}_{\mathrm{next}}$ resident in registers, thereby consuming $B_r \times B_c \times \text{sizeof}(\text{float})$ more register space for each threadblock. This heightened register consumption can limit the extent to which one can select larger block dimensions—a commonly used performance strategy that is itself register-intensive. Ultimately, careful profiling should be used to choose the most effective configuration in light of these trade-offs.
\paragraph{3-stage pipelining} By extending the preceding 2-stage approach, we introduce a 3-stage pipeline variant, enabling further concurrency between the second WGMMA invocation and the softmax calculation. Although this strategy may further improve Tensor Core occupancy, it amplifies register demands, as yet another stage’s data must be buffered, increasing the complexity of balancing tile dimensions and pipeline depth. A comprehensive treatment of the 3-stage approach and experimental analysis is given in ~\cref{sec:3-stage}.

\subsection{Low-precision with FP8}
\label{sec:algofp8}

\textbf{Efficiency: layout transformations.} Performing the forward pass of \textsc{FlashAttention-3} using FP8 precision introduces further complications regarding layout compatibility that are not present when using FP16.

To begin, it is important to highlight that the input tensors $\mathbf{Q}$, $\mathbf{K}$, and $\mathbf{V}$ are generally stored such that the head dimension is contiguous. However, in order to comply with the k-major memory layout requirement of FP8 WGMMA for the second GEMM, it is necessary for $\mathbf{V}$—specifically, its tiles as placed in SMEM—to have contiguity along the sequence length dimension. Since TMA loading does not alter the axis along which the data is contiguous, we are compelled to either (1) perform a pre-processing transpose of $\mathbf{V}$ within GMEM, or (2) carry out a tile-wise in-kernel transpose within SMEM after data is loaded. For approach (1), this could entail (1a) incorporating the transpose operation into the epilogue phase of a preceding layer such as rotary embedding, or (1b) launching a dedicated pre-processing kernel for the transpose\footnote{A well-optimized transpose kernel can approach the device’s maximum bandwidth~\citep{colfax_cutlass_transpose_2024}.} to swap the order of the sequence length and head strides. Nevertheless, approach (1a) presents significant challenges for integration in generic library code, while (1b) is highly inefficient when memory bandwidth is a limiting factor, as it frequently is during inference.

For FP8 \textsc{FlashAttention-3}, we select option (2) as our approach. The in-kernel transpose leverages the LDSM (\verb|ldmatrix|) and STSM (\verb|stmatrix|) instructions, which enable a warp of threads to collaboratively transfer data between SMEM and RMEM at 128-byte blocks.\footnote{PTX documentation specifies that LDSM/STSM operations copy $8 \times 8$ matrices where each element is 16 bits \cite[\S 9.7.13.4.15-16]{ptx}; for FP8, we achieve compatibility by packing pairs of 8-bit values, enabling the use of LDSM/STSM at lower precision. However, the transpose variants of LDSM/STSM are unable to split these packed 8-bit values, which imposes additional register operations between LDSM and STSM to complete a tile-wise transpose; full details are elided here.} These instructions are highly efficient in terms of register usage, enabling their execution directly within the producer warpgroup, and are capable of handling both memory copy and layout transposition tasks. In addition, beyond the first iteration, we can schedule the transpose operation for the upcoming $\mathbf{V}$ tile concurrently, so that it overlaps with the execution of the two WGMMAs that consume the prior $\mathbf{V}$ and the current $\mathbf{K}$ tiles.

Second, our analysis indicates that, in contrast to FP16, the memory organization of the FP32 accumulator in an FP8 WGMMA differs from the storage arrangement assumed for operand A when it resides in registers. \cref{fig:rmem_accum} and \cref{fig:rmem_operand} illustrate partial representations of both layouts, showing how register entries are assigned per thread in the sequence indicated. Utilizing byte permute instructions, we are able to rearrange the output of the initial WGMMA's accumulator, reformatting it to be compatible both with the requirements of a subsequent WGMMA invocation and with the organization of the $\mathbf{V}$ tile created by the in-kernel transpose. Concretely, according to \cref{fig:rmem_accum}, the order is changed as follows:
$$\{ \verb|d0 d1 d4 d5 d2 d3 d6 d7| \},$$
and this register permutation repeats for every 8-byte segment. This transformation acts as a column permutation of the $\mathbf{P}$ tile in logical terms, meaning, for instance, that columns $0189$ are brought to the front. For the second WGMMA to produce the correct output arrangement, we can coordinate by applying a corresponding row permutation when transposing the $\mathbf{V}$ tile within the kernel.\footnote{By leveraging this flexibility of an in-kernel transpose, we avoid the need for shuffle instructions, which would otherwise be required to redistribute register values among threads, as discussed previously in~\citep{colfax_fp8_flashattention_2024}.}

\textbf{Accuracy: block quantization and incoherent processing.} The FP8 (e4m3) format allocates only 3 bits for the mantissa and 4 bits for the exponent, leading to greater numerical imprecision relative to FP16 or BF16 formats. In addition, large-scale models frequently display outlier values~\citep{dettmers2208llm, sun2024massive}—entries with much larger magnitudes compared to the majority, which presents challenges for quantization. Typically, per-tensor scaling~\citep{micikevicius2022fp8} is implemented, where each tensor (for example, $\mathbf{Q}$, $\mathbf{K}$, and $\mathbf{V}$) is associated with its own scalar. To address numerical inaccuracies arising in attention computations with FP8, we adopt two main strategies:

\begin{enumerate}

\item \textbf{Block quantization}: In this method, a single scalar is maintained for each block; specifically, for each of $\mathbf{Q}$, $\mathbf{K}$, and $\mathbf{V}$, the tensor is divided into blocks of dimension $B_r \times d$ or $B_c \times d$. Each block is quantized independently. This quantization step can be seamlessly combined with preceding operations such as rotary embedding, without incurring any additional computational overhead, as these operations are predominantly constrained by memory bandwidth. Since the \textsc{FlashAttention-3} algorithm is inherently designed to operate on blocks, we are able to appropriately rescale each block of $\mathbf{S}$ to compensate for block-wise quantization, again with no extra computation required.

\item \textbf{Incoherent processing}: To mitigate the effects of outlier values, prior to quantization to FP8, we pre-multiply both $\mathbf{Q}$ and $\mathbf{K}$ by a random orthogonal matrix $\mathbf{M}$. Because $\mathbf{M}$ satisfies $\mathbf{M} \mathbf{M}^\top = I$, the product $(\mathbf{Q} \mathbf{M})(\mathbf{K} \mathbf{M})^\top$ yields the same result as $\mathbf{Q} \mathbf{K}^\top$, ensuring that this transformation does not alter the final attention calculation. The random transformation effectively disperses outlier effects since each element in $\mathbf{Q} \mathbf{M}$ or $\mathbf{K} \mathbf{M}$ becomes a randomized combination of the original tensor elements, decreasing the quantization error. Our implementation follows the strategies of \citet{chee2024quip} and~\citet{tseng2024quip}, constructing $\mathbf{M}$ as a product of random diagonal matrices with entries of $\pm 1$ and a Hadamard matrix. This enables the transformation to be performed in $O(d \log d)$ time, much faster than the $O(d^2)$ cost of general matrix multiplication, and it can be integrated with the rotary embedding step without adding to computational expense.
\end{enumerate}

Our experiments demonstrate that these two methods can collectively reduce numerical error by up to 2.6$\times$, as shown in \cref{sec:numerical_error}.

\section{Empirical Validation}
\label{sec:exp}

Our implementation of \textsc{FlashAttention-3} leverages WGMMA and TMA abstractions, as provided by CUTLASS~\citep{Thakkar_CUTLASS_2023}, to assess both its performance and precision.

\begin{itemize}

\item \textbf{Benchmarking attention.} We evaluate the execution time of \textsc{FlashAttention-3} for various sequence lengths, conducting a comparison against a conventional PyTorch implementation, \textsc{FlashAttention-2}, a version of \textsc{FlashAttention-2} in Triton that leverages H100-specific instructions, as well as a vendor-tuned \textsc{FlashAttention-2} from cuDNN, optimized specifically for H100 GPUs. Our experiments demonstrate that \textsc{FlashAttention-3} delivers up to 2.0$\times$ speedup relative to \textsc{FlashAttention-2} and achieves a 1.5$\times$ increase in performance over the Triton-based \textsc{FlashAttention-2}. Furthermore, \textsc{FlashAttention-3} attains throughput as high as 740 TFLOPs/s, which corresponds to 75\% of the H100 GPU's peak theoretical TFLOPs/s.
\item \textbf{Ablation study.} We provide evidence that the inclusion of warp-specialization and the GEMM-softmax pipelining strategy in our methodology substantially drives the observed performance improvements in \textsc{FlashAttention-3}.
\item \textbf{Accuracy of FP8 attention.} Our results show that the application of block quantization along with incoherent processing mitigates numerical errors in FP8 \textsc{FlashAttention-3}, reducing error by a factor of 2.6$\times$.
\end{itemize}

\subsection{Benchmarking Attention}
\label{subsec:benchmark_attn}

We evaluate the execution time of several attention mechanisms on an H100 80GB SXM5 GPU under varying configurations (both with and without a causal mask, head sizes of 64 or 128), using FP16 input precision. Performance outcomes are presented in~\cref{fig:benchmark_attn_fwd} and~\cref{fig:benchmark_attn_bwd}. The results indicate that \textsc{FlashAttention-3} achieves speedups of approximately 1.5-2.0$\times$ relative to \textsc{FlashAttention-2} during the forward pass, and is about 1.5-1.75$\times$ faster in the backward pass. When compared to a conventional attention baseline, \textsc{FlashAttention-3} can deliver improvements ranging from 3$\times$ to 16$\times$. Notably, for sequence lengths at or beyond 1k, \textsc{FlashAttention-3} outperforms the cuDNN vendor library (closed source), which is highly optimized for H100 GPUs.

\paragraph{Benchmark settings:} We vary the sequence length as 512, 1k, ..., 16k, and set batch size so that the total number of tokens is 16k. We set the hidden dimension to 2048, and head dimension to be either 64, 128, or 256 (i.e., 32 heads, 16 heads, or 8 heads). To calculate the FLOPs of the forward pass, we use:

\begin{equation*}
  4 \cdot \text{seqlen}^2 \cdot \text{head dimension} \cdot \text{number of heads}.
\end{equation*}

By employing causal masking, the total is halved since, on average, only about 50% of the elements are processed. For the backward pass, the number of FLOPs is estimated by scaling the forward pass FLOPs by a factor of 2.5, reflecting the presence of two matrix multiplications during the forward computation and five during the backward computation because of recomputation requirements.

Additionally, we evaluate the forward pass runtime using FP8 under comparable conditions. The outcomes for headdim 256 are depicted in ~\cref{fig:benchmark_attn_fp8}, while a comprehensive set of results can be found in \cref{sec:benchmark_attn_fp8_full}.

\subsection{Ablation Study: 2-Stage Pipelining Experiments}
\label{sec:2-stage-pipelining-experiments}

We conduct an ablation study of both the 2-stage WGMMA-softmax pipelining and warp-specialization in the context of non-causal FP16 \textsc{FlashAttention-3}, using the fixed parameter set $\{\text{batch}, \text{seqlen}, \text{nheads}, \text{hdim}\} = \{ 4, 8448, 16, 128\}$. As shown in~\cref{table:ablation_pipelining}, our enhancements—incorporating asynchrony via warp-specialization and overlapping the GEMM and softmax operations—result in a marked performance improvement, boosting throughput from 570 to 661 TFLOPs.

\subsection{Numerical Error Validation}
\label{sec:numerical_error}

Given the recent attention to the numerical error~\citep{golden2024flash} in \textsc{FlashAttention}, we evaluate the performance of \textsc{FlashAttention-2}, \textsc{FlashAttention-3}, and a conventional attention implementation relative to a reference baseline computed in FP64. To mimic outlier features and activations commonly encountered in LLMs~\citep{dettmers2208llm, sun2024massive}, we sample the elements of $\mathbf{Q}, \mathbf{K}, \mathbf{V}$ according to the following distribution:

\begin{equation*}
  \mathcal{N}(0, 1) + \mathcal{N}(0, 100) \cdot \mathrm{Bernoulli}(0.001).
\end{equation*}

Specifically, every element follows a normal distribution with a mean of zero and a standard deviation of 1; however, for 0.1\% of the elements, an additional independent term is included, itself sampled from a normal distribution with a standard deviation of 10. The root mean squared error (RMSE) results are reported in~\cref{table:numerical_error}. In FP16 format, both \textsc{FlashAttention-2} and \textsc{FlashAttention-3} obtain RMSE values that are 1.7$\times$ lower than those produced by the conventional implementation, as they maintain intermediate (softmax) results in FP32. In the case of FP8, the baseline attention uses per-tensor scaling, carries out matrix multiplications in FP32, and stores intermediate softmax results in FP16. Due to the employment of block quantization and incoherent processing, \textsc{FlashAttention-3} in FP8 surpasses this baseline in accuracy by a factor of 2.6$\times$.

\section{Dicussion, Limitations, Conclusion}
\label{sec:discussion}

Through \textsc{FlashAttention-3}, we illustrate that leveraging modern programming strategies and hardware innovations—specifically asynchrony and low-precision computation—can significantly enhance both the speed and accuracy of attention mechanisms. Our approach achieves a 1.5-2.0$\times$ acceleration in attention computation over \textsc{FlashAttention-2}, while also cutting FP8 numerical error by a factor of 2.6 compared to traditional per-tensor quantization. There are, however, some current limitations that remain to be addressed: optimizing specifically for LLM inference scenarios, incorporating a persistent kernel framework into the FP8 kernel,\footnote{In our benchmarks, FP16 \textsc{FlashAttention-3} benefits from a persistent kernel design and advanced load balancing, whereas FP8 \textsc{FlashAttention-3} does not, which partially accounts for its comparatively lower performance in small sequence lengths and causal masking when compared to the FP8 cuDNN kernels.} and further investigating the impact of low-precision attention during large-scale training regimes. Although the present study targets Hopper GPUs, we anticipate that these methods are broadly transferable to other types of hardware accelerators. Ultimately, we envision that advancements in attention primitives with improved speed and accuracy could enable new possibilities, particularly for applications requiring extended context lengths.

\end{document}